\chapter{结果与分析}

\section{职业阶段与媒介重心}

示例分析将佐佐木希的跨媒介轨迹概括为模特起步、影视拓展、形象调整与平台再建构四个阶段。表~\ref{tab:jlau-stage} 展示不同阶段的媒介重心和主要形象功能。

\begin{table}[htbp]
  \centering
  \caption{佐佐木希跨媒介职业阶段的示例比较}
  \label{tab:jlau-stage}
  \begin{tabularx}{\textwidth}{>{\raggedright\arraybackslash}p{2.5cm}>{\raggedright\arraybackslash}p{2.8cm}X}
    \toprule
    阶段 & 媒介重心 & 示例性形象特征 \\
    \midrule
    模特起步 & 平面杂志与广告 & 以稳定视觉风格建立初始识别度 \\
    影视拓展 & 电视剧、电影与综艺 & 借助角色和节目互动扩展公众认知 \\
    形象调整 & 访谈与品牌活动 & 在职业表达与公共讨论之间重组叙事 \\
    平台再建构 & 社交平台与生活方式内容 & 通过连续、日常化更新增强亲近感 \\
    \bottomrule
  \end{tabularx}
\end{table}

\section{媒介可见度的示例变化}

图~\ref{fig:jlau-visibility} 使用模拟指数演示折线图的标题、图例和坐标布局。数值不代表真实传播测量结果。

\begin{figure}[htbp]
  \centering
  \begin{tikzpicture}[x=2.45cm,y=0.062cm]
    \draw[->,thick] (0,0) -- (4.45,0) node[right] {阶段};
    \draw[->,thick] (0,0) -- (0,112) node[above] {示例指数};
    \foreach \y in {20,40,60,80,100}{
      \draw[black!15] (0,\y) -- (4.15,\y);
      \node[left] at (0,\y) {\small \y};
    }
    \foreach \x/\label in {0.55/模特起步,1.55/影视拓展,2.55/形象调整,3.55/平台再建构}{
      \node[below,align=center] at (\x,-3) {\small \label};
    }
    \draw[very thick,black!80] plot coordinates {(0.55,58) (1.55,86) (2.55,52) (3.55,78)};
    \foreach \x/\y in {0.55/58,1.55/86,2.55/52,3.55/78}{
      \fill[black!80] (\x,\y) circle[radius=1.6pt];
      \node[above] at (\x,\y) {\small \y};
    }
  \end{tikzpicture}
  \caption{不同职业阶段媒介可见度的模拟变化}
  \label{fig:jlau-visibility}
\end{figure}

\section{讨论}

示例结果提示，跨媒介转型并非从一种平台简单迁移到另一种平台。杂志建立的视觉资本、影视角色提供的叙事空间和社交平台形成的更新节奏会相互叠加。对佐佐木希这类公开人物进行研究时，应避免把阶段变化归因于单一事件，而要比较媒介结构、职业选择和受众互动的共同影响。

\section{小结}

本章通过一幅趋势图和一张阶段比较表展示了常见结果章节的排版方式。使用者可以替换坐标、数据和文字，同时保留题注位置、章节编号与交叉引用结构。
